\chapter{Authorization}
\label{ch:authorization}

\section{Product authorization}

\textbf{Status:} NOT APPLICABLE.

Expo Snips implements no authorization model. There are no roles, permissions,
policies, access-control lists, middleware guards or protected resources in the
product. Relevant evidence:

\begin{itemize}
  \item The CLI has no notion of a user or a role; all subcommands are available to
        anyone who can execute the binary.
  \item The extension has no \texttt{when} activation clauses beyond command and
        language activation, and no user-dependent behaviour.
  \item The registry is public, read-only static hosting; every recipe and every
        template is readable by anyone who can reach the host.
  \item There is no role/permission data model, no RLS (the product has no
        database) and no policy engine. All \textbf{NOT APPLICABLE}.
\end{itemize}

\section{Effective privilege boundaries}

Although there is no authorization \emph{model}, the system does have privilege
boundaries that matter operationally:

\begin{longtable}{@{}L{4.4cm}L{5.4cm}L{4.6cm}@{}}
\caption{Effective privilege boundaries}\\
\toprule
\textbf{Boundary} & \textbf{What is possible} & \textbf{What is not} \\
\midrule
\endfirsthead
\toprule
\textbf{Boundary} & \textbf{What is possible} & \textbf{What is not} \\
\midrule
\endhead
Local file writes & Generate files at any path declared by a recipe, relative to the
project root. & Read arbitrary files (the tool only reads \texttt{package.json},
\texttt{tsconfig.json}, the manifest and files it is about to overwrite). \\
Child process execution & Run the detected package manager; run \texttt{npx
create-expo-app} during \texttt{init}. & Run arbitrary commands from recipe data ---
recipe \texttt{dependencies} strings are inserted into the install command line. \\
Network access & Fetch JSON from the configured registry base. & Upload data;
transmit credentials; authenticate. \\
Registry content & Any published recipe is trusted and rendered verbatim. & Verify
provenance or signature. \\
\bottomrule
\end{longtable}

\section{A UI access consideration worth recording}

The catalog contains three recipes that are \emph{not} end-user screens but shared
UI primitives (\texttt{ui/button}, \texttt{ui/input}, \texttt{ui/utils}). They are
listed by \texttt{expo-snips list} and are selectable from the CLI picker and the
VS Code Quick Pick exactly like application screens. Nothing in the CLI prevents a
user from running:

\begin{lstlisting}[language=bash]
expo-snips add ui/button
\end{lstlisting}

which writes \texttt{components/ui/button.tsx} (and, because \texttt{ui/button}
declares \texttt{primitives: ["ui/utils"]}, also \texttt{lib/utils.ts}) directly
into the project, with \texttt{styling} forced to \texttt{nativewind} and
\texttt{ui} forced to \texttt{shadcn} by the primitive-resolution logic.

\textbf{Assessment (INFERRED).} This is not a security problem, but it is a
usability and presentation concern: internal primitives are indistinguishable from
application screens in every discovery surface. A future revision could filter
\texttt{ui/*} ids from the user-facing pickers or mark them as internal.

\section{Recommendation}

If a hosted, curated registry is ever offered, authorization will become relevant
in at least two places:

\begin{enumerate}
  \item \textbf{Write access} to publish recipes --- currently anyone with write
        access to the repository (and therefore the CI token) can publish, and
        there is no review gate beyond pull-request review.
  \item \textbf{Read access} to private/company catalogs --- the configurable
        registry has no support for authentication headers or tokens, so a private
        registry must currently be reachable without credentials (for example, on a
        trusted network or behind a proxy).
\end{enumerate}

\textbf{Status:} RECOMMENDED. Neither behaviour exists today.
